%\documentclass{article}
%\usepackage{graphicx} % Required for inserting images

%\title{Hojjat-AAAI-2027}
%\author{morteza.rezanejad }
%\date{July 2026}

%\begin{document}

%\maketitle

%\section{Introduction}

%\end{document}
% ============================================================
%  AAAI 2027 Anonymous Submission
%  "Weakly-Supervised 3D Anomaly Segmentation via
%   Soft Mask-Coupled Diffusion Inpainting"
% ============================================================
\documentclass[letterpaper]{article} % DO NOT CHANGE THIS
\usepackage[submission]{aaai2027}    % DO NOT CHANGE THIS
\usepackage[hyphens]{url}            % DO NOT CHANGE THIS
\usepackage{graphicx}                % DO NOT CHANGE THIS
\urlstyle{rm}                        % DO NOT CHANGE THIS
\def\UrlFont{\rm}                    % DO NOT CHANGE THIS
\usepackage{natbib}                  % DO NOT CHANGE THIS
\usepackage{caption}                 % DO NOT CHANGE THIS
\frenchspacing                       % DO NOT CHANGE THIS

% Recommended (algorithms)
\usepackage{algorithm}
\usepackage{algorithmic}

% Recommended (tables)
\usepackage{booktabs}

% Additional math / formatting (permitted)
\usepackage{amsmath}
\usepackage{amssymb}
\usepackage{mathtools}
\usepackage{bm}
\usepackage{xcolor}
\usepackage{multirow}

% ── Convenience macros ────────────────────────────────────────
\newcommand{\mt}{m_t}
\newcommand{\xhat}{\hat{x}_0}
\newcommand{\rt}{r_t}
\newcommand{\phealthy}{p_\phi(y\!=\!\text{healthy}\mid x_t,t)}
\newcommand{\todo}[1]{\textcolor{red}{\textbf{[TODO: #1]}}}

\pdfinfo{/TemplateVersion (2027.1)}

\setcounter{secnumdepth}{1}

% ── Title / Author (anonymous) ────────────────────────────────
\title{Weakly-Supervised 3D Anomaly Segmentation via\\
       Soft Mask-Coupled Diffusion Inpainting}

\author{
    Anonymous Submission
}
\affiliations{
    Anonymous Institution
}

% =============================================================
\begin{document}
\maketitle

% =============================================================
\begin{abstract}
Weakly supervised anomaly segmentation in volumetric medical imaging remains challenging because dense lesion annotations are costly and difficult to obtain at scale, whereas coarse clinical cues such as slice-level bounding boxes are far more accessible in routine practice. We propose a 3D diffusion inpainting framework that learns a generative model of healthy anatomy from nodule-free chest CT and localizes abnormalities at inference by identifying regions that cannot be plausibly reconstructed as normal tissue. Our method trains a volumetric diffusion model exclusively on healthy patches, then applies RePaint-style inpainting conditioned on weak per-slice bounding-box prompts to synthesize plausible normal tissue within candidate regions; the resulting residual between the observed scan and the inpainted reconstruction defines the anomaly mask. An iterative refinement procedure further sharpens localization by updating the inpainting region where large reconstruction errors persist, enabling voxel-level segmentation from limited supervision without using lesion masks during training. Evaluated on a fixed LIDC-IDRI hold-out benchmark against contemporary foundation and supervised segmentation approaches, our method outperforms most existing foundation models in native 3D and remains competitive with strong 2D alternatives despite substantially weaker prompts. A central advantage of this formulation is that the model captures the distribution of normal anatomical variation rather than lesion-specific appearance, making it naturally suited to cross-dataset and cross-organ generalization; we additionally analyze distributional differences between synthetic and real lesion embeddings to characterize domain shift and support broader deployment. Together, these results demonstrate that generative diffusion inpainting offers a scalable, principled, and effective paradigm for weakly supervised anomaly segmentation in medical imaging.
\end{abstract}

% =============================================================
\section{Introduction}
\label{sec:intro}

Volumetric anomaly segmentation in 3D medical imaging is a
prerequisite for radiomics analysis, malignancy stratification,
and longitudinal monitoring. Fully supervised 3D segmentation
models achieve strong performance on LIDC-IDRI, the only large-
scale lung CT dataset with voxel-level mask annotations
(${\approx}1{,}000$ cases). The vast majority of clinical 3D
imaging datasets---including the National Lung Screening Trial
(NLST, 26,254 subjects) and analogous repositories across liver,
kidney, and brain---carry only cheap \emph{weak labels}: bounding
boxes, centroids, or RECIST diameters. This annotation bottleneck
is universal across organ types and imaging modalities.

State-of-the-art interactive foundation models
\citep{nninteractive2025,segvol2024,sam_med3d2023,medsam2_2025}
accept box/point prompts but segment ``the thing in the box'' by
boundary contrast: they have no mechanism to distinguish
pathological from normal tissue. This fails critically for
sub-6\,mm nodules and ground-glass opacities (GGOs), where
intensity boundaries are absent or ambiguous across all
modalities.

We reframe weakly supervised anomaly segmentation as
\emph{healthy-tissue counterfactual generation} via 3D diffusion
inpainting. Our model is trained only on healthy patches; at
inference, given a weak label it inpaints the indicated region as
plausible healthy tissue, and the residual localizes the anomaly.
A novel \emph{soft mask-coupled reverse process} refines this
residual into a tight 3D pseudo-segmentation mask entirely inside
the generative sampling dynamics. Because the model learns only
``normal,'' it is anomaly-agnostic and transfers across organs and
modalities without retraining on lesion data.

\paragraph{Contributions.}
Our contributions span three groups:

\begin{itemize}
\item \textbf{Diffusion \& Sampling (C1--C4).}
    Soft Gaussian mask initialization from weak labels~(C1) and a
    soft mask-coupled reverse diffusion kernel~(C2), in which the
    anomaly mask $m_t$ evolves as a latent variable inside the
    reverse process---the first such formulation in
    reconstruction-based anomaly detection.
    A multi-timestep discrepancy score~(C3) exploiting the
    lesion-size-dependent information content of the per-step
    residual $r_t$, and mask-scoped classifier guidance~(C4) that
    actively suppresses pathological reconstructions within the
    evolving mask.

\item \textbf{Training (C5--C6).}
    A four-component synthetic anomaly augmentation suite~(C5)
    and an anomaly-aware hinge regularizer~(C6) that sharpen the
    healthy-tissue prior and improve robustness to domain shift,
    theoretically coupled to C2 and C4 through the soft mask.

\item \textbf{Generalisation (C7--C9).}
    Cross-dataset generalisation on LIDC-IDRI, NLST, LUNA16~(C7);
    cross-organ transfer from lung to liver/kidney CT~(C8); and
    cross-modality transfer from CT to brain MRI~(C9)---
    demonstrating the framework is not lung- or CT-specific.
\end{itemize}

% =============================================================
\section{Related Work}
\label{sec:related}

\paragraph{Supervised 3D Lesion Segmentation.}
nnU-Net \citep{nnunet2021} is the de facto upper bound for
supervised CT segmentation, auto-configuring a 3D U-Net to any
dataset. Transformer-based architectures (SwinUNETR, nnFormer)
improve on LIDC-IDRI but require full mask supervision and degrade
sharply outside their training distribution, especially on
sub-6\,mm nodules due to extreme class imbalance.

\paragraph{Weakly Supervised Lesion Segmentation.}
DeepEM \citep{deepem2018} trains 3D CNNs from location labels via
EM over latent boxes. RECIST-based supervision
\citep{zhou2023recist} converts 2D diameters into 3D pseudo-
labels. Most relevantly, \citet{rectflow2026} use a 3D latent
diffusion model's generative guidance to suppress nodule signal in
reconstruction on LUNA16---the closest existing work. Key
differences: their model is not trained on healthy-only data;
their mask is not evolved inside the reverse process; they provide
no classifier guidance, regularizer, or cross-organ/modality
evaluation.

\paragraph{Reconstruction-Based Anomaly Detection.}
AnoDDPM \citep{anoddpm2022} introduced partial diffusion with
simplex noise for brain MRI, training on healthy data and
computing reconstruction error as the anomaly score
($+25.5\%$ Dice over f-AnoGAN). Extensions include latent-space
partial diffusion \citep{latent_ad2023}, Synomaly noise for
ultrasound \citep{synomaly2024}, counterfactual generation for
brain tumours \citep{fontanella2023}, IterMask2
\citep{itermask2024}, and FDP \citep{fdp2026} for brain MRI.
All are 2D or 2D-stacked, use a fixed corruption region, and do
\emph{not} evolve the mask inside the reverse process.

\paragraph{Adaptive Mask Inpainting.}
AMI-Net \citep{aminet2024} proposes an adaptive mask generator for
industrial anomaly detection (MVTec AD, BTAD). MAD-paint
\citep{madpaint2025} introduces mask-uncertainty metrics for
diffusion inpainting. Neither is 3D, medical, nor couples the mask
inside the reverse diffusion kernel.

\paragraph{Foundation Models for Medical Segmentation.}
MedSAM \citep{medsam2024} fine-tunes SAM on medical data but
remains 2D. SAM-Med3D \citep{sam_med3d2023} retrains SAM natively
in 3D. SegVol \citep{segvol2024} is a universal interactive CT
segmenter. MedSAM2 \citep{medsam2_2025} applies SAM2's video
memory to 3D volumes. nnInteractive \citep{nninteractive2025} won
the CVPR 2025 Foundation Models for Interactive 3D Biomedical
Image Segmentation Challenge. All are boundary-detection paradigms that lack an explicit
anomaly-versus-normal generative prior; they also tend to degrade
on small or low-contrast lesions.

% =============================================================
\section{Methodology}
\label{sec:method}
%==============================================================================

We formulate lung-nodule localization as \emph{anomaly-aware generative
inpainting}.
A 3D generative model is trained exclusively on \emph{healthy} lung tissue,
thereby encoding a prior over normal CT appearance.
At test time, a candidate region is erased and reconstructed under that prior;
voxels that cannot be explained by healthy tissue leave a large reconstruction
residual.
By iteratively refining the inpaint region from this residual, we obtain a
binary nodule mask without training a dedicated segmentation network.

\paragraph{Notation.}
Let \(\mathbf{x}\in\mathbb{R}^{D\times H\times W}\) denote a CT patch after HU
windowing and normalization (Sec.~\ref{sec:method:data}).
Let \(\mathbf{y}\in\{0,1\}^{D\times H\times W}\) be the ground-truth nodule mask
(used only to construct the \emph{initial} prompt and for evaluation).
We write \(\mathbf{m}\in\{0,1\}^{D\times H\times W}\) for the \emph{keep} mask
(\(\mathbf{m}=1\) known~/~observed, \(\mathbf{m}=0\) inpaint hole) and
\(\mathbf{h}=1-\mathbf{m}\) for the hole.
The network \(v_\theta\) predicts a rectified-flow velocity field.

%------------------------------------------------------------------------------
\subsection{Problem Formulation}
\label{sec:method:problem}
%------------------------------------------------------------------------------

Given a thoracic CT patch containing a candidate nodule, the goal is to
recover a binary segmentation \(\hat{\mathbf{y}}\) of the lesion.
Conventional discriminative segmenters map \(\mathbf{x}\mapsto\hat{\mathbf{y}}\)
directly.
In contrast, we leverage the observation that nodules are \emph{deviations from
healthy lung anatomy}: if a generative model of healthy tissue is forced to
reconstruct an erased nodule region, the reconstruction \(\hat{\mathbf{x}}\)
disagrees with \(\mathbf{x}\) precisely where anomalous tissue lies.
The final hole \(\mathbf{h}^\star\) obtained after residual-driven refinement
is taken as the prediction,
\begin{equation}
\hat{\mathbf{y}} \;=\; \mathbf{h}^\star.
\end{equation}

%------------------------------------------------------------------------------
\subsection{Data and Preprocessing}
\label{sec:method:data}
%------------------------------------------------------------------------------

We use LIDC-IDRI chest CT volumes~\cite{armato2011lidc}.
Patches of size \(64^{3}\) are extracted at resampling to spacing
\((0.7,0.7,1.5)\,\mathrm{mm}\) and retained only when the lung mask occupies at
least \(25\%\) of the patch.
Intensities are clipped to \([-1000,400]\,\mathrm{HU}\) and linearly mapped to
\([0,1]\):
\begin{equation}
\tilde{x}
=
\mathrm{clip}\!\Biggl(
\frac{x - \mathrm{HU}_{\min}}{\mathrm{HU}_{\max}-\mathrm{HU}_{\min}},\,0,1
\Biggr),
\quad\\
\mathrm{HU}_{\min}=-1000,\;
\mathrm{HU}_{\max}=400.
\end{equation}
Training uses \emph{nodule-free} (healthy) patches; inference and evaluation use
\emph{positive} nodule patches with LIDC consensus annotations.
Patient-level disjoint splits prevent leakage across train~/~validation~/~test.

%------------------------------------------------------------------------------
\subsection{Healthy Tissue Prior via Rectified Flow}
\label{sec:method:rflow}
%------------------------------------------------------------------------------

\paragraph{Generative model.}
We train an unconditional 3D U-Net \(v_\theta\)~\cite{ronneberger2015unet}
to transport noise to healthy CT under
\emph{rectified flow}~\cite{liu2023flow,lipman2023flowmatching}.
Unlike score-based DDPM training, rectified flow learns a nearly straight
probability path between \(\boldsymbol{\varepsilon}\sim\mathcal{N}(\mathbf{0},\mathbf{I})\)
and data \(\mathbf{x}_0\), yielding efficient few-step ODE sampling.

\paragraph{Training path.}
Let \(T\) denote the continuous time horizon (we use \(T{=}1000\) with continuous
timesteps and timestep transform~\cite{esser2024sd3}).
A sample is formed by linear interpolation
\begin{equation}
\mathbf{x}_t
=
\Bigl(1-\frac{t}{T}\Bigr)\mathbf{x}_0
+
\frac{t}{T}\,\boldsymbol{\varepsilon},
\qquad
t\in[0,T],
\end{equation}
and the regression target is the constant velocity along this path,
\begin{equation}
\mathbf{u}(\mathbf{x}_t,t)
=
\mathbf{x}_0 - \boldsymbol{\varepsilon}.
\end{equation}
The network is optimized with an \(\ell_1\) velocity objective
\begin{equation}
\mathcal{L}(\theta)
=
\mathbb{E}_{\mathbf{x}_0,\boldsymbol{\varepsilon},t}
\Biggl[
\bigl\|
v_\theta(\mathbf{x}_t,t)
-
(\mathbf{x}_0-\boldsymbol{\varepsilon})
\bigr\|_1
\Biggr].
\end{equation}
Volumes are standardized by a per-batch scale factor
\(s=1/\mathrm{std}(\mathbf{x}_0)\) before diffusion, and rescaled at the end of
sampling.

\paragraph{Architecture and optimization.}
The U-Net has four resolution levels with channels
\([64,128,128,128]\), two residual blocks per level, and attention at the
bottleneck (\texttt{num\_head\_channels}\({=}16\)).
Inputs and outputs are single-channel.
We use Adam with learning rate \(10^{-4}\), gradient clipping at \(1.0\),
bfloat16 mixed precision, and medium spatial~/~intensity augmentations
(random flips; intensity shift \(\pm 0.05\) and scale \(\pm 0.10\) in
normalized space).
No masks or ControlNet residuals are used during training: the prior is a pure
healthy-tissue generator.

%------------------------------------------------------------------------------
\subsection{Mask-Guided Rectified-Flow Inpainting}
\label{sec:method:repaint}
%------------------------------------------------------------------------------

At inference we perform RePaint-style~\cite{lugmayr2022repaint} known-region
injection under the rectified-flow ODE, without jump resampling.

Given keep-mask \(\mathbf{m}\) and observation \(\mathbf{x}_0\), sampling starts
from \(\mathbf{x}\sim\mathcal{N}(\mathbf{0},\mathbf{I})\) in scaled space and
follows an \(N\)-step Euler discretization of the learned velocity
(\(N{=}40\) in our default setting).
At each time \(t\), known voxels are overwritten by a noised copy of the
observation while unknown voxels evolve freely:
\begin{equation}
\label{eq:inject}
\mathbf{x}
\;\leftarrow\;
\mathbf{m}\odot \tilde{\mathbf{x}}_t(\mathbf{x}_0)
\;+\;
(1-\mathbf{m})\odot \mathbf{x},
\end{equation}
where \(\tilde{\mathbf{x}}_t(\mathbf{x}_0)\) denotes the rectified-flow
forward perturbation of \(\mathbf{x}_0\) at time \(t\).
A network step then updates
\begin{equation}
\mathbf{x}
\;\leftarrow\;
\mathbf{x}
+
v_\theta(\mathbf{x},t)\,\Delta t,
\qquad
\Delta t=\frac{t-t_{\mathrm{next}}}{T}.
\end{equation}
After \(N\) steps and unscaling, we obtain an inpainted volume
\(\hat{\mathbf{x}}\) that is consistent with \(\mathbf{x}_0\) outside the hole
and synthesized from the healthy prior inside the hole.

%------------------------------------------------------------------------------
\subsection{Initial Hole from Slice-wise Boxes}
\label{sec:method:init}
%------------------------------------------------------------------------------

Prompted methods require a spatial query.
Following standard interactive segmentation practice, we construct an
\emph{initial} hole from per-slice axis-aligned boxes of the nodule annotation
(or, at deployment, of a coarse detector / user scribbles).
For each axial slice \(z\) containing foreground,
\begin{equation}
\mathbf{h}^{(0)}_z
=
\mathrm{AABB}_{2\mathrm{D}}\!\bigl(\mathbf{y}_z\bigr)
\oplus\, p,
\end{equation}
with padding \(p{=}1\) voxel, stacked over \(z\) into a 3D binary hole
\(\mathbf{h}^{(0)}\).
This deliberately over-covers the lesion and provides a conservative starting
region for refinement.

%------------------------------------------------------------------------------
\subsection{Iterative Residual Mask Refinement}
\label{sec:method:refine}
%------------------------------------------------------------------------------

A fixed initial hole is imperfect: it may include healthy soft tissue or miss
irregular lesion boundaries.
We therefore iterate inpainting and residual thresholding
(Algorithm~\ref{alg:refine}).

After inpainting with hole \(\mathbf{h}^{(r)}\), define the absolute residual
in normalized intensity space,
\begin{equation}
\label{eq:residual}
\mathbf{d}^{(r)}
=
\bigl|\mathbf{x}_0 - \hat{\mathbf{x}}^{(r)}\bigr|.
\end{equation}
Large residuals indicate tissue that the healthy prior cannot explain.
The next hole is obtained by thresholding inside a local search band around
the previous hole:
\begin{equation}
\label{eq:update}
\mathbf{h}^{(r+1)}
=
\mathbb{1}\!\bigl[\mathbf{d}^{(r)} > \tau\bigr]
\;\odot\;
\mathrm{Dilate}\!\bigl(\mathbf{h}^{(r)};\,\rho\bigr),
\end{equation}
with residual threshold \(\tau{=}0.05\) and 3D dilation radius \(\rho{=}3\).
If the resulting mask contains fewer than \(\eta{=}10\) voxels, we retain
\(\mathbf{h}^{(r)}\) (stability safeguard).
We run \(R\) refinement rounds (\(R{=}2\) in our best configuration;
\(R{=}3\) in ablations), performing \(R{+}1\) inpaint passes in total.
No morphological open~/~close post-processing is applied in the default method.

\begin{algorithm}[t]
\caption{Iterative residual mask refinement for nodule localization}
\label{alg:refine}
\begin{algorithmic}[1]
\Require CT patch \(\mathbf{x}_0\), initial hole \(\mathbf{h}^{(0)}\),
         RF inpainter \(\mathcal{G}_\theta\), rounds \(R\), threshold \(\tau\),
         dilation \(\rho\), min.\ voxels \(\eta\)
\Ensure Predicted nodule mask \(\hat{\mathbf{y}}\)
\For{\(r = 0,\dots,R\)}
    \State \(\hat{\mathbf{x}}^{(r)}
            \gets \mathcal{G}_\theta\!\bigl(\mathbf{x}_0,\,\mathbf{m}{=}1{-}\mathbf{h}^{(r)}\bigr)
            \)
            \Comment{Eq.~\eqref{eq:inject}}
    \If{\(r = R\)}
        \State \textbf{break}
    \EndIf
    \State \(\mathbf{d}^{(r)} \gets |\mathbf{x}_0 - \hat{\mathbf{x}}^{(r)}|\)
    \State \(\mathbf{s} \gets \mathrm{Dilate}(\mathbf{h}^{(r)};\,\rho)\)
    \State \(\mathbf{h}' \gets \mathbb{1}[\mathbf{d}^{(r)} > \tau]\odot\mathbf{s}\)
    \State \(\mathbf{h}^{(r+1)} \gets
            \mathbf{h}'\) if \(|\mathbf{h}'|_0 \ge \eta\), else \(\mathbf{h}^{(r)}\)
\EndFor
\State \Return \(\hat{\mathbf{y}} \gets \mathbf{h}^{(R)}\)
\end{algorithmic}
\end{algorithm}

The refinement loop couples geometry and appearance: boxes supply a spatial
prior, while residuals correct the mask toward the true lesion extent.
Because sampling is conditioned on known healthy context through
Eq.~\eqref{eq:inject}, each iterate becomes progressively more selective.

%------------------------------------------------------------------------------
\subsection{Relation to Discriminative and Promptable Segmenters}
\label{sec:method:discuss}
%------------------------------------------------------------------------------

Discriminative segmenters and foundation models (e.g., SAM-Med2D) map prompts
to masks in a single feed-forward pass.
Our method instead \emph{tests} a spatial hypothesis against a generative
healthy prior and updates the hypothesis via reconstruction error.
The same initial boxes used to prompt interactive models initialize our hole;
refinement then specializes the region without requiring paired nodule-dense
training of a segmenter.
This makes the approach complementary to promptable models and attractive when
healthy tissue is abundant but curated lesion masks are scarce.

%------------------------------------------------------------------------------
\subsection{Implementation Details}
\label{sec:method:impl}
%------------------------------------------------------------------------------

Unless noted otherwise we use \(N{=}40\) RF inference steps, \(R{=}2\)
refinement rounds, \(\tau{=}0.05\), \(\rho{=}3\), and \(p{=}1\).
Checkpoints are selected by validation velocity loss on held-out healthy
patches.
All experiments share the same HU normalization and \(64^{3}\) geometry as
training.
We report Dice, IoU, precision, and recall between \(\hat{\mathbf{y}}\) and
LIDC nodule masks on a held-out positive-patch test set.

Ablations over \(N\in\{30,40,60,80\}\) and \(R\in\{2,3\}\) show that a moderate
sampling budget with \(R{=}2\) yields the strongest Dice on the full test cache
(1658 cases), reaching a mean Dice of approximately \(0.734\), competitive with
SAM-Med2D (\(0.737\)) under matched box prompts.

%==============================================================================
\section*{Default Hyperparameters}
%==============================================================================

\begin{table}[h]
\centering
\caption{Default hyperparameters used in our method.}
\label{tab:hparams}
\begin{tabular}{@{}ll@{}}
\toprule
\textbf{Setting} & \textbf{Value} \\
\midrule
Patch size & \(64\times 64\times 64\) \\
Spacing & \(0.7\times 0.7\times 1.5\,\mathrm{mm}\) \\
HU window & \([-1000,\,400]\) \\
Scheduler & Rectified flow (continuous) \\
Train timesteps \(T\) & \(1000\) \\
U-Net channels & \([64,128,128,128]\) \\
Loss & \(\ell_1\) velocity regression \\
Learning rate & \(10^{-4}\) \\
Inference steps \(N\) & \(40\) \\
Refine rounds \(R\) & \(2\) \\
Residual threshold \(\tau\) & \(0.05\) \\
Search dilation \(\rho\) & \(3\) \\
Box padding \(p\) & \(1\) voxel \\
Morphology & off \\
\bottomrule
\end{tabular}
\end{table}
% =============================================================
\section{Experiments}
\label{sec:experiments}

\subsection{Datasets}
\textbf{Lung CT.}
LIDC-IDRI \citep{lidc2011} (1,018 scans, four-radiologist
consensus voxel masks; \emph{oracle evaluation}).
NLST \citep{nlst2011} (26,254 subjects, box/centroid weak labels;
\emph{scale test}).
LUNA16 \citep{luna2017} (\emph{FROC benchmark}).

\noindent\textbf{Cross-organ CT (C8).}
MSD Task~03 liver tumours \citep{msd2019} or KiTS kidney tumours
\citep{kits2023} (full mask GT; box-only weak labels at
inference).

\noindent\textbf{Cross-modality MRI (C9).}
BraTS \citep{brats2015} brain tumours (full mask GT; box-only weak
labels at inference).

\subsection{Baselines}
We compare against:
(i)~nnU-Net 3D \citep{nnunet2021} (supervised upper bound);
(ii)~MedSegDiff-V2 \citep{medsegdiff2024} (supervised diffusion
segmentation, AAAI~2024);
(iii)~AnoDDPM reimplemented natively in 3D \citep{anoddpm2022}
(unsupervised reconstruction, no weak labels);
(iv)~VAE-based 3D anomaly detection (non-diffusion reconstruction
baseline);
(v)~AMI-Net adapted to 3D CT \citep{aminet2024};
(vi)~MedSAM \citep{medsam2024} (2D, box-prompted, slice-by-slice);
(vii)~SAM-Med3D \citep{sam_med3d2023};
(viii)~SegVol \citep{segvol2024};
(ix)~MedSAM2 \citep{medsam2_2025};
(x)~nnInteractive \citep{nninteractive2025} (all 3D, box-prompted;
nnInteractive is the CVPR~2025 challenge winner).

\subsection{Main Results}

\todo{Table~\ref{tab:main}: Dice/IoU on LIDC-IDRI and FROC on
LUNA16 across all baselines. Fill after experiments complete.}

\begin{table}[t]
\centering
\caption{Main segmentation results on LIDC-IDRI (Dice/IoU, full
mask GT) and LUNA16 (FROC). Best weakly supervised result in
\textbf{bold}; best overall \underline{underlined}.}
\label{tab:main}
\begin{tabular}{lccc}
\toprule
Method & Supervision & LIDC Dice & LUNA FROC \\
\midrule
nnU-Net 3D          & Full mask     & \todo{} & \todo{} \\
MedSegDiff-V2       & Full mask     & \todo{} & \todo{} \\
\midrule
AnoDDPM-3D          & Unsupervised  & \todo{} & \todo{} \\
VAE-3D AD           & Unsupervised  & \todo{} & \todo{} \\
\midrule
MedSAM (2D)         & Box prompt    & \todo{} & \todo{} \\
SAM-Med3D           & Box prompt    & \todo{} & \todo{} \\
SegVol              & Box prompt    & \todo{} & \todo{} \\
MedSAM2             & Box prompt    & \todo{} & \todo{} \\
nnInteractive       & Box prompt    & \todo{} & \todo{} \\
\midrule
\textbf{Ours (Full)}& Box/centroid  & \todo{} & \todo{} \\
\bottomrule
\end{tabular}
\end{table}

\paragraph{Stratified Analysis.}
Table~\ref{tab:strat} reports Dice stratified by nodule diameter,
demonstrating our advantage in the sub-6\,mm regime against all
baselines including nnInteractive---the primary empirical claim of
this paper.

\begin{table}[t]
\centering
\caption{Dice stratified by nodule diameter on LIDC-IDRI.}
\label{tab:strat}
\begin{tabular}{lccc}
\toprule
Method & $<6$\,mm & $6$--$10$\,mm & $>10$\,mm \\
\midrule
nnU-Net 3D    & \todo{} & \todo{} & \todo{} \\
nnInteractive & \todo{} & \todo{} & \todo{} \\
SegVol        & \todo{} & \todo{} & \todo{} \\
\textbf{Ours} & \todo{} & \todo{} & \todo{} \\
\bottomrule
\end{tabular}
\end{table}

\subsection{Simulated Weak-Label Experiment}
To directly validate our method in the annotation regime it
targets, we strip LIDC-IDRI masks to bounding boxes/centroids and
run the full pipeline using \emph{only} these degraded labels,
then evaluate against the real withheld masks. This addresses the
primary reviewer concern: ``does the method work when you actually
only have weak labels?''

\subsection{Ablation Study}

Table~\ref{tab:ablation} ablates each contribution in sequence,
isolating the effect of each component.

\begin{table}[t]
\centering
\caption{Ablation on LIDC-IDRI Dice (LIDC) and sub-6\,mm Dice
(Small).}
\label{tab:ablation}
\begin{tabular}{lcc}
\toprule
Configuration                              & LIDC Dice & Small \\
\midrule
(0) Hard-box inpainting (baseline)         & \todo{} & \todo{} \\
(1) + Soft Gaussian init (C1)              & \todo{} & \todo{} \\
(2) + Mask-coupled reverse (C2)            & \todo{} & \todo{} \\
(3) + Multi-timestep score (C3)            & \todo{} & \todo{} \\
(4) + Classifier guidance (C4)             & \todo{} & \todo{} \\
(5) + Synthetic augmentation (C5)          & \todo{} & \todo{} \\
(6) + Hinge regularizer (C6) = Full model  & \todo{} & \todo{} \\
\bottomrule
\end{tabular}
\end{table}

Additional ablations cover: initial $\sigma$ sensitivity;
$\lambda$/$\tau_t$ annealing schedule; guidance scale $s$; DDIM
step count; and each augmentation component independently.

\subsection{Cross-Organ and Cross-Modality Generalisation}

Table~\ref{tab:generalization} reports MSD liver Dice (C8) and
BraTS Dice (C9) compared against AnoDDPM-3D, IterMask2, and
FDP-AAAI2026 in the box-prompted weak-label regime. We also
report downstream malignancy classification AUC from radiomic
features extracted from NLST pseudo-masks, compared against
features from LIDC-IDRI expert masks, directly validating the
``weakly supervised segmentation for downstream radiomics''
framing.

\begin{table}[t]
\centering
\caption{Generalisation across organs (liver CT) and modalities
(brain MRI), box-prompted weak-label setting.}
\label{tab:generalization}
\begin{tabular}{lccc}
\toprule
Method & Liver Dice (C8) & BraTS Dice (C9) & NLST AUC \\
\midrule
AnoDDPM-3D    & \todo{} & \todo{} & --- \\
IterMask2     & ---     & \todo{} & --- \\
FDP           & ---     & \todo{} & --- \\
nnInteractive & \todo{} & \todo{} & --- \\
\textbf{Ours} & \todo{} & \todo{} & \todo{} \\
\bottomrule
\end{tabular}
\end{table}

% =============================================================
\section{Scope, Limitations, and Conclusion}
\label{sec:conclusion}

\subsection{Scope}
Within lung CT, the framework is anomaly-agnostic: solid nodules,
part-solid nodules, GGOs, and any localized deviation from healthy
parenchyma are detected via the same residual mechanism. Cross-
organ and cross-modality scope is empirically demonstrated (C8,
C9). Extension to ultrasound and PET requires domain-specific
healthy-prior training and is reserved for future work. The
framework is applicable to any clinical dataset with bounding box
or centroid annotations and a healthy-tissue corpus for training.

\subsection{Limitations}
\begin{itemize}
    \item \textbf{Inference speed.} Iterative reverse diffusion
    with per-step mask updates and classifier guidance is slower
    than a single network forward pass. DDIM acceleration is
    applied but per-nodule latency exceeds nnU-Net/nnInteractive.
    Consistency-model distillation is planned future work.

    \item \textbf{Classifier dependency.} C4's classifier is
    trained on LIDC-IDRI, introducing a semi-supervised element;
    it must be retrained per new organ/modality for C8/C9.

    \item \textbf{Patch size.} The $64^3$ patch-level framework
    may fragment very large or diffuse lesions at boundaries;
    multi-scale hierarchical inference is future work.

    \item \textbf{Modality scope.} Cross-modality validation
    covers brain MRI only; ultrasound and PET are not evaluated
    in this paper.
\end{itemize}

\subsection{Conclusion}
We presented a unified framework for weakly supervised 3D anomaly
segmentation requiring no voxel-level annotations at any stage.
Our soft mask-coupled reverse diffusion (C1, C2), multi-timestep
discrepancy score and classifier guidance (C3, C4), synthetic
augmentation and hinge regularizer (C5, C6), and cross-dataset/-
organ/-modality generalisation (C7--C9) constitute a coherent
system coupled through the evolving soft mask $m_t$. Because the
model is trained only on healthy tissue and never on any specific
pathology, it is anomaly-agnostic by construction: the same
pipeline detects solid nodules, GGOs, liver lesions, and brain
tumours using only weak labels---without retraining on lesion
data. The framework outperforms interactive foundation models
specifically in the sub-6\,mm, low-contrast regime where
boundary-detection paradigms fail, and generalises across organs
and modalities by design.

% =============================================================
\section*{Acknowledgments}
[Omitted for anonymous review.]

% =============================================================
\bibliography{paper}

\end{document}
